% from ega2-1-fr.tex
\oldpage[II]{109}

avec l'isomorphisme $\psi':U\to U'$, que
$q_1\circ\psi=u_k\circ q_2\circ\psi$. Mais comme $q_1\circ\psi=j$,
$q_2\circ\psi=\varphi$, notre assertion résulte de la commutativité du
diagramme (\hyperref[II.5.6.1.1-fr]{5.6.1.1}), compte tenu de la définition de
$u_k$.

E) Il est clair que comme $U$, et par suite $U'$, est irréductible, il en est
de même de $X'$ dans la construction précédente, et le morphisme $f$ est donc
birationnel (\hyperref[I.2.2.9-fr]{I, 2.2.9}). Si en outre $X$ est réduit, il
en est de même de $U'$, donc $X'$ est réduit
(\hyperref[I.9.5.9-fr]{I, 9.5.9}). Cela achève la démonstration.

\begin{corollary}[5.6.2]
\label{II.5.6.2-fr}
On suppose vérifiée l'une des hypothèses a), b) de
(\hyperref[II.5.6.1-fr]{5.6.1}). Pour que $X$ soit propre sur $S$, il faut et
il suffit qu'il existe un schéma $X'$ projectif sur $S$ et un $S$-morphisme
surjectif $f:X'\to X$ (qui est alors projectif par
(\hyperref[II.5.5.5-fr]{5.5.5, (v)})). Lorsqu'il en est ainsi, on peut en
outre choisir $f$ de sorte qu'il existe un ouvert dense $U$ de $X$ tel que la
restriction de $f$ à $f^{-1}(U)$ soit un isomorphisme
$f^{-1}(U)\simeq U$ et que $f^{-1}(U)$ soit dense dans $X'$. Si en outre $X$
est irréductible (resp. réduit) on peut supposer qu'il en est de même de
$X'$~; lorsque $X$ et $X'$ sont irréductibles, $f$ est un morphisme
birationnel.
\end{corollary}

La condition est suffisante en vertu de (\hyperref[II.5.5.3-fr]{5.5.3}) et
(\hyperref[II.5.4.3-fr]{5.4.3, (ii)}). Elle est nécessaire, car avec les
notations de (\hyperref[II.5.6.1-fr]{5.6.1}), si $X$ est propre sur $S$, $X'$
est propre sur $S$, car il est projectif sur $X$, donc propre sur $X$
(\hyperref[II.5.5.3-fr]{5.5.3}), et notre assertion résulte de
(\hyperref[II.5.4.2-fr]{5.4.2, (ii)})~; en outre, comme $X'$ est
quasi-projectif sur $S$, il est projectif sur $S$ en vertu de
(\hyperref[II.5.5.3-fr]{5.5.3}).

\begin{corollary}[5.6.3]
\label{II.5.6.3-fr}
Soit $S$ un préschéma localement noethérien, $X$ un $S$-schéma de type fini
sur $S$, $f_0:X\to S$ son morphisme structural. Pour que $X$ soit propre sur
$S$, il faut et il suffit que pour tout morphisme de type fini $S'\to S$,
$(f_0)_{(S')}:X_{(S')}\to S'$ soit un morphisme fermé. Il suffit même que
cette condition soit vérifiée pour tout $S$-préschéma de la forme
$S'=S\otimes_{\mathbf Z}\mathbf Z[T_1,\ldots,T_n]$ ($T_i$ indéterminées).
\end{corollary}

La condition étant évidemment nécessaire, montrons qu'elle est suffisante.
La question étant locale sur $S$ et $S'$ (\hyperref[II.5.4.1-fr]{5.4.1}), on
peut supposer $S$ et $S'$ affines et noethériens. En vertu du lemme de Chow,
il existe un $S$-schéma projectif $P$, une immersion $j:X'\to P$ et un
morphisme projectif et surjectif $f:X'\to X$, tels que le diagramme
\[
\xymatrix{
X\ar[d]_{f_0} & X'\ar[l]_f\ar[d]^j\\
S & P\ar[l]_r
}
\]
soit commutatif. Comme $P$ est de type fini sur $S$, la première hypothèse
entraîne que la projection $q_2:X\times_SP\to P$ est un morphisme fermé. Mais
l'immersion $j$ est composée de $q_2$ et du morphisme $f\times1$ de
$X'\times_SP$ dans $X\times_SP$~; or $f$, étant projectif, est propre
(\hyperref[II.5.5.3-fr]{5.5.3}), donc $f\times1$ est fermé. On en conclut que
$j$ est une immersion fermée, et par suite est propre
(\hyperref[II.5.4.2-fr]{5.4.2, (i)}). En outre le morphisme structural
$r:P\to S$ est projectif, donc propre (\hyperref[II.5.5.3-fr]{5.5.3}), donc
$f_0\circ f=r\circ j$ est propre
(\hyperref[II.5.4.2-fr]{5.4.2, (ii)})~; finalement, comme $f$ est surjectif,
$f_0$ est propre en vertu de (\hyperref[II.5.4.3-fr]{5.4.3}).

Pour prouver la proposition en utilisant seulement la seconde hypothèse, il
suffit de montrer qu'elle entraîne la première. Or, si $S'$ est affine et de
type fini sur $S=\operatorname{Spec}(A)$,

